\chapter{自然语言处理：预训练}

词是离散符号。预训练把词或词元映到向量，使共现与上下文进入内积。
先写 word2vec 与 GloVe，再写子词与 BERT 的掩蔽语言建模、下一句预测。

\section{词嵌入（word2vec）}

词向量 $\bm{v}\in\mathbb{R}^{d}$ 是词的特征。词嵌入是映射 $w\mapsto\bm{v}_w$。

\subsection{为何独热向量是一个糟糕的选择}

词表 $\mathcal{V}$ 上的独热 $\bm{e}_w$ 两两正交，余弦
\[
\frac{\bm{x}^{\mathsf T}\bm{y}}{\|\bm{x}\|\|\bm{y}\|}\in[-1,1]
\]
对任何不同词都是 $0$，无法表达相似。

\subsection{自监督的word2vec}

跳元（skip-gram）用中心词预测上下文；连续词袋（CBOW）用上下文预测中心词。
二者都用语料中的共现，不另需标签。

\subsection{跳元模型（Skip-Gram）}

由 “loves” 生成两侧词的联合，并假定条件独立：
\begin{align*}
&P(\text{the},\text{man},\text{his},\text{son}\mid\text{loves})
\\
&\approx
P(\text{the}\mid\text{loves})\,P(\text{man}\mid\text{loves})
\,P(\text{his}\mid\text{loves})\,P(\text{son}\mid\text{loves}).
\end{align*}
一般地，中心词向量 $\bm{v}_c$、上下文向量 $\bm{u}_o$，
\[
P(w_o\mid w_c)
=
\frac{\exp(\bm{u}_o^{\mathsf T}\bm{v}_c)}{\sum_{i\in\mathcal{V}}\exp(\bm{u}_i^{\mathsf T}\bm{v}_c)}.
\]
长度为 $T$、窗口 $m$ 的语料似然
\[
\prod_{t=1}^{T}
\prod_{\substack{-m\leq j\leq m\\ j\neq 0}}
P\bigl(w^{(t+j)}\mid w^{(t)}\bigr).
\]

\subsubsection{训练}

最小化
\[
-\sum_{t=1}^{T}
\sum_{\substack{-m\leq j\leq m\\ j\neq 0}}
\log P\bigl(w^{(t+j)}\mid w^{(t)}\bigr).
\]
对数条件概率
\[
\log P(w_o\mid w_c)
=
\bm{u}_o^{\mathsf T}\bm{v}_c
-\log\sum_{i\in\mathcal{V}}\exp(\bm{u}_i^{\mathsf T}\bm{v}_c).
\]
对 $\bm{v}_c$ 的梯度
\begin{align*}
\frac{\partial\log P(w_o\mid w_c)}{\partial\bm{v}_c}
&=
\bm{u}_o
-\sum_{j\in\mathcal{V}}P(w_j\mid w_c)\,\bm{u}_j.
\end{align*}
每步与 $|\mathcal{V}|$ 成正比。

\subsection{连续词袋（CBOW）模型}

\[
P(\text{loves}\mid\text{the},\text{man},\text{his},\text{son}).
\]
上下文平均 $\bar{\bm{v}}_o=\frac1{2m}(\bm{v}_{o_1}+\cdots+\bm{v}_{o_{2m}})$，
\begin{align*}
P(w_c\mid\mathcal{W}_o)
&=
\frac{\exp(\bm{u}_c^{\mathsf T}\bar{\bm{v}}_o)}
{\sum_{i\in\mathcal{V}}\exp(\bm{u}_i^{\mathsf T}\bar{\bm{v}}_o)}.
\end{align*}
序列似然
\[
\prod_{t=1}^{T}
P\bigl(w^{(t)}\mid w^{(t-m)},\ldots,w^{(t-1)},w^{(t+1)},\ldots,w^{(t+m)}\bigr).
\]

\subsubsection{训练}

\[
-\sum_{t=1}^{T}
\log P\bigl(w^{(t)}\mid w^{(t-m)},\ldots,w^{(t+m)}\bigr).
\]
\[
\log P(w_c\mid\mathcal{W}_o)
=
\bm{u}_c^{\mathsf T}\bar{\bm{v}}_o
-\log\sum_{i\in\mathcal{V}}\exp(\bm{u}_i^{\mathsf T}\bar{\bm{v}}_o).
\]
\[
\frac{\partial\log P(w_c\mid\mathcal{W}_o)}{\partial\bm{v}_{o_i}}
=
\frac1{2m}
\Biggl(
\bm{u}_c-\sum_{j\in\mathcal{V}}P(w_j\mid\mathcal{W}_o)\,\bm{u}_j
\Biggr).
\]

\subsection{小结}

跳元：$P(w_o\mid w_c)\propto\exp(\bm{u}_o^{\mathsf T}\bm{v}_c)$。
CBOW：$P(w_c\mid\mathcal{W}_o)\propto\exp(\bm{u}_c^{\mathsf T}\bar{\bm{v}}_o)$。
Softmax 分母扫过整个 $\mathcal{V}$。

\subsection{练习}

核验：梯度对 $|\mathcal{V}|$ 的复杂度，以及 $\bm{u}^{\mathsf T}\bm{v}$ 与余弦的关系。

\section{近似训练}

精确 Softmax 每步求和 $|\mathcal{V}|$ 项。负采样与层序 Softmax 把代价降到常数或对数。

\subsection{负采样}

把“中心–上下文是否共现”当成二分类
\[
P(D=1\mid w_c,w_o)=\sigma(\bm{u}_o^{\mathsf T}\bm{v}_c),\qquad
\sigma(x)=\frac1{1+\exp(-x)}.
\]
正对似然再乘 $K$ 个噪声词的负例
\begin{align*}
P(w^{(t+j)}\mid w^{(t)})
&=
P(D=1\mid w^{(t)},w^{(t+j)})
\prod_{k=1}^{K}
P(D=0\mid w^{(t)},w_k).
\end{align*}
损失
\begin{align*}
-\log P(w^{(t+j)}\mid w^{(t)})
&=
-\log\sigma(\bm{u}_{i_{t+j}}^{\mathsf T}\bm{v}_{i_t})
-\sum_{k=1}^{K}
\log\sigma(-\bm{u}_{h_k}^{\mathsf T}\bm{v}_{i_t}).
\end{align*}
每步与 $K$ 线性，与 $|\mathcal{V}|$ 无关。

\subsection{层序Softmax}

词表做成二叉树。从根到叶 $w_o$ 的路径长度为 $L(w_o)$，
\begin{align*}
P(w_o\mid w_c)
&=
\prod_{j=1}^{L(w_o)-1}
\sigma\Bigl(
[\![n(w_o,j+1)=\mathrm{leftChild}(n(w_o,j))]\!]
\,\bm{u}_{n(w_o,j)}^{\mathsf T}\bm{v}_c
\Bigr).
\end{align*}
路径上向左与向右对应 $+\bm{u}^{\mathsf T}\bm{v}$ 与 $-\bm{u}^{\mathsf T}\bm{v}$。
对任意 $w_c$ 有 $\sum_{w\in\mathcal{V}}P(w\mid w_c)=1$。
代价为 $O(\log|\mathcal{V}|)$。

\subsection{小结}

负采样用 $1+K$ 次 $\sigma(\bm{u}^{\mathsf T}\bm{v})$；
层序 Softmax 用根到叶的 $O(\log|\mathcal{V}|)$ 次。

\subsection{练习}

核验：$\sum_w P(w\mid w_c)=1$ 在二叉树上是否成立，以及噪声分布 $P(w)$ 如何取。

\section{用于预训练词嵌入的数据集}

把 PTB 文本变成跳元 $+$ 负采样所需的小批量。

\subsection{读取数据集}

一行一句、空格分词。频次低于阈值的词并入未知词元。

\subsection{下采样}

词 $w_i$ 被丢弃的概率
\[
P(w_i)=\max\Bigl(1-\sqrt{\frac{t}{f(w_i)}},0\Bigr),
\]
$f$ 为相对频率，$t$ 为阈值。高频词更常被丢掉。

\subsection{中心词和上下文词的提取}

对每个中心位置均匀抽窗口半径 $1$ 到 $m$，距离不超过该半径的词为上下文。

\subsection{负采样}

噪声 $w$ 的抽样概率正比于 $f(w)^{0.75}$。每对正例抽 $K$ 个噪声词。

\subsection{小批量加载训练实例}

第 $i$ 个样本有 $n_i$ 个上下文与 $m_i$ 个噪声。
拼接到长度 $\max_i(n_i+m_i)$ 并补零；掩码标出非填充；标签区分正负例。

\subsection{整合代码}

迭代器一次产出中心词、上下文加噪声、掩码与标签。

\subsection{小结}

下采样降高频；批量用掩码对齐变长的正负例。

\subsection{练习}

核验：关闭下采样对每个 epoch 词元数与墙钟时间的影响。

\section{预训练word2vec}

在 PTB 上用负采样训练跳元。

\subsection{跳元模型}

两个嵌入：中心与上下文（含噪声）。

\subsubsection{嵌入层}

权重 $\bm{E}\in\mathbb{R}^{|\mathcal{V}|\times d}$。索引 $i$ 取第 $i$ 行。
批量索引形状 $(B,L)$ 时输出 $(B,L,d)$。

\subsubsection{定义前向传播}

中心 $(B,1,d)$ 与上下文–噪声 $(B,\mathrm{max\_len},d)$ 做批量点积，
得 $(B,1,\mathrm{max\_len})$，每个元素是 $\bm{v}_c^{\mathsf T}\bm{u}$。

\subsection{训练}

损失只在非填充位置上平均。

\subsubsection{二元交叉熵损失}

对每个正负位置用
$-\bigl[y\log\sigma(s)+(1-y)\log(1-\sigma(s))\bigr]$，
再按掩码平均，与负采样公式一致。

\subsubsection{初始化模型参数}

两张嵌入表，行数为 $|\mathcal{V}|$，例如 $d=100$。

\subsubsection{定义训练阶段代码}

每个批量：前向点积、掩码交叉熵、反传更新两张表。

\subsection{应用词嵌入}

查询词 $\bm{v}$，在词表中取余弦最大的 $k$ 个 $\bm{v}'$。

\subsection{小结}

跳元前向是 $\bm{v}_c^{\mathsf T}\bm{u}$；训练是掩码二元交叉熵。
下游用余弦找近邻。

\subsection{练习}

核验：同一中心词在不同轮次重抽上下文与噪声，对更新方差的影响。

\section{全局向量的词嵌入（GloVe）}

用全局共现计数 $x_{ij}$ 解释跳元，再用平方损失拟合 $\log x_{ij}$。

\subsection{带全局语料统计的跳元模型}

\[
q_{ij}
=
\frac{\exp(\bm{u}_j^{\mathsf T}\bm{v}_i)}{\sum_{k\in\mathcal{V}}\exp(\bm{u}_k^{\mathsf T}\bm{v}_i)},
\]
共现加权交叉熵
\[
-\sum_{i\in\mathcal{V}}\sum_{j\in\mathcal{V}}x_{ij}\log q_{ij}
=
-\sum_{i\in\mathcal{V}}x_i\sum_{j\in\mathcal{V}}p_{ij}\log q_{ij},
\]
其中 $x_i=\sum_j x_{ij}$，$p_{ij}=x_{ij}/x_i$。

\subsection{GloVe模型}

\[
\sum_{i\in\mathcal{V}}\sum_{j\in\mathcal{V}}
h(x_{ij})
\bigl(\bm{u}_j^{\mathsf T}\bm{v}_i+b_i+c_j-\log x_{ij}\bigr)^2.
\]
$h$ 对大计数截断，避免极高频对支配损失。中心向量与上下文向量在此目标下地位对称。

\subsection{从条件概率比值理解GloVe模型}

希望
\[
f(\bm{u}_j,\bm{u}_k,\bm{v}_i)\approx\frac{p_{ij}}{p_{ik}}.
\]
取指数内积比
\[
\frac{\exp(\bm{u}_j^{\mathsf T}\bm{v}_i)}{\exp(\bm{u}_k^{\mathsf T}\bm{v}_i)}
\approx
\frac{p_{ij}}{p_{ik}},
\]
从而
\[
\bm{u}_j^{\mathsf T}\bm{v}_i+b_i+c_j\approx\log x_{ij}.
\]

\subsection{小结}

GloVe 用加权平方损失拟合 $\log x_{ij}$，而不是对整个 $\mathcal{V}$ 做 Softmax。
比值 $p_{ij}/p_{ik}$ 由内积差解释。

\subsection{练习}

核验：用窗口内距离重加权 $p_{ij}$ 后目标如何变，以及 $b_i$ 与 $c_i$ 是否可对换。

\section{子词嵌入}

变形、复合词共享字符片段。把词向量写成子词向量之和。

\subsection{fastText模型}

词 $w$ 的字符 $n$-gram 集合 $\mathcal{G}_w$（含词本身），
\[
\bm{v}_w=\sum_{g\in\mathcal{G}_w}\bm{z}_g.
\]
跳元仍用 $\bm{u}_o^{\mathsf T}\bm{v}_c$，但 $\bm{v}_c$ 可加性分解。未见词只要子词在词表中即可表示。

\subsection{字节对编码（Byte Pair Encoding）}

从单字符开始，反复合并当前语料中最频的相邻符号对，得到可变长子词，词表大小预先固定。
不跨词边界合并。

\subsection{小结}

fastText：$\bm{v}_w=\sum_{g\in\mathcal{G}_w}\bm{z}_g$。
BPE 贪心合并最频对，得到固定词表上的可变长子词。

\subsection{练习}

核验：从 $n$ 个初始符号得到词表大小 $m$ 需要多少次合并。

\section{词的相似性和类比任务}

大语料上预训练的向量可直接用于近邻与类比。

\subsection{加载预训练词向量}

GloVe 有 $50/100/300$ 维；fastText 有 $300$ 维英文向量。
词表含四十万词加未知词元。

\subsection{应用预训练词向量}

相似度用余弦；$k$ 近邻在词表上穷举。

\subsubsection{词相似度}

对查询 $w$，返回 $\arg\max_{w'\neq w}\cos(\bm{v}_w,\bm{v}_{w'})$ 的前 $k$ 个。

\subsubsection{词类比}

$a:b::c:d$ 中求
\[
d=\operatorname*{argmin}_{w}\bigl\|\bm{v}_w-(\bm{v}_c+\bm{v}_b-\bm{v}_a)\bigr\|
\]
（或最大余弦）。例如性别与首都–国家。

\subsection{小结}

近邻是 $\cos(\bm{v},\bm{v}')$；类比是 $\bm{v}_c+\bm{v}_b-\bm{v}_a$ 的最近邻。

\subsection{练习}

核验：词表很大时如何用近似近邻代替全表扫描。

\section{来自Transformers的双向编码器表示（BERT）}

word2vec/GloVe 的 $\bm{v}_w$ 与上下文无关。BERT 用双向 Transformer 编码器，使表示依赖整句。

\subsection{从上下文无关到上下文敏感}

上下文无关：$f(x)$ 只摄入词元 $x$。
多义词在不同句中应不同；$f(x,\text{上下文})$ 才能区分。

\subsection{从特定于任务到不可知任务}

ELMo 双向但下游仍要定制结构；GPT 任务无关但从左到右。
微调时 GPT 更新全部解码器参数，ELMo 常冻结。

\subsection{BERT：把两个最好的结合起来}

双向编码 $+$ 下游只加一个浅输出层，并微调全部编码器参数。

\subsection{输入表示}

单句：$\langle\mathrm{cls}\rangle+\,\text{词元}+\,\langle\mathrm{sep}\rangle$。
句对：$\langle\mathrm{cls}\rangle+A+\langle\mathrm{sep}\rangle+B+\langle\mathrm{sep}\rangle$。
输入嵌入为词元嵌入、片段嵌入与位置嵌入之和。

\subsection{预训练任务}

编码器输出每个词元（含特殊词元）的向量，供两个自监督头使用。

\subsubsection{掩蔽语言模型（Masked Language Modeling）}

随机选 $15\%$ 词元做预测。对被选词元：以 $80\%$ 换成 $\langle\mathrm{mask}\rangle$，
$10\%$ 换成随机词，$10\%$ 保持原词，再用双向上下文预测原词。
避免微调阶段从未见过 $\langle\mathrm{mask}\rangle$ 导致的不一致。

\subsubsection{下一句预测（Next Sentence Prediction）}

句对一半为真实邻句（正），一半第二句随机（负）。
$\langle\mathrm{cls}\rangle$ 的表示经单隐层 MLP 做二分类。

\subsection{整合代码}

总损失为 MLM 与 NSP 损失的和。前向返回编码、MLM logits 与 NSP logits。

\subsection{小结}

BERT 输入是三类嵌入之和。预训练
\[
\ell=\ell_{\mathrm{MLM}}+\ell_{\mathrm{NSP}},
\]
MLM 给双向上下文，NSP 给句对关系。

\subsection{练习}

核验：MLM 相对从左到右 LM 需要更多还是更少步才能收敛，以及 GELU 与 ReLU 的差别。

\section{用于预训练BERT的数据集}

在 WikiText-2 上生成 MLM 与 NSP 样本。保留标点、大小写与数字。

\subsection{为预训练任务定义辅助函数}

从段落造句对，再对 BERT 序列做掩蔽。

\subsubsection{生成下一句预测任务的数据}

以 $1/2$ 概率取下句为真邻句，否则从语料随机抽一句。截断使总长不超过 $\mathrm{max\_len}$。

\subsubsection{生成遮蔽语言模型任务的数据}

在非特殊词元中抽约 $15\%$ 位置；按 $80/10/10$ 替换；返回改写后的词元、预测位置与原词标签。

\subsection{将文本转换为预训练数据集}

填充到固定长，附片段下标、有效长度掩码、MLM 标签与 NSP 标签。
一条样本对应一对句子上的两个任务。

\subsection{小结}

WikiText-2 样本是填充后的
$(\text{词元},\text{片段},\text{MLM 位置与标签},\text{NSP 标签})$。

\subsection{练习}

核验：不用句号切句、不过滤稀有词元时词表有多大。

\section{预训练BERT}

用上一节样本训练小型 BERT，再用编码器表示文本。

\subsection{预训练BERT}

原论文 $\mathrm{BERT_{BASE}}$：$12$ 层、$768$ 隐、$12$ 头，约 $1.1\times 10^8$ 参数；
$\mathrm{BERT_{LARGE}}$：$24$ 层、$1024$、$16$ 头，约 $3.4\times 10^8$。
演示用 $2$ 层、$128$ 隐、$2$ 头。
批量损失为该批量 MLM 与 NSP 损失之和。

\subsection{用BERT表示文本}

单句或句对加上特殊词元后编码。
$\langle\mathrm{cls}\rangle$ 位是整段表示；同一词元（如 crane）在不同句中的向量不同。

\subsection{小结}

预训练后 $f_{\mathrm{BERT}}(x,\text{上下文})$ 随上下文变。
$\langle\mathrm{cls}\rangle$ 与各词元向量都可作下游输入。

\subsection{练习}

核验：为何实验中 MLM 损失高于 NSP，以及 $\mathrm{max\_len}=512$ 配大模型时显存是否溢出。
